% ---
% id: meat-core/thm:choi
% kind: thm
% title: "The Choi program"
% summary: >-
%   In trusted-source mode the program is exactly a jointly convex
%   minimisation over sector Choi blocks, with one coupling constraint per
%   Alice sector and affine or conic rows. Every support frame gives the value
%   h_low, since only surjectivity of the parametrisation enters it.
% prominence: standard
% status: implemented
% uses: [lem:choi-marginal, def:support-frame, def:feasible, lem:admissible-alice-subspace, def:keyactive, thm:convexity, lem:choi-factor, lem:joint-sectors, def:joint-sectors, def:declaration, prop:exact-cancellation, foundations/def:sector-map]
% assumes: []
% refs: []
% provenance: newly-derived
% implements: certiqs_key.meat_core.objective.choi:choi_parametrization
% lean:
% ---

Applied block by block along the sectors, the parametrisation of
Lemma~\ref{lem:choi-marginal} rewrites the program of
Definition~\ref{def:feasible} in Choi blocks, the variables a solver works
in, without changing its value on any support frame.

\begin{theorem}\label{thm:choi}
Let the declaration be in \tref{trusted-source mode}, fix a
\tref{support frame} (Definition~\ref{def:support-frame}), apply
Lemma~\ref{lem:choi-marginal} with $B'=\Qsys$, and restrict to
sector-block-diagonal operators. Then the assignment
$J=\bigoplus_{s}J_s\mapsto\omega=\chi_{\Qsys}(J)$, with
$\omega_s=\chi_{\Qsys_s}(J_s)\coloneqq
(W_{\tau}\otimes\idon{\Qsys_s})J_s(W_{\tau}^{\dagger}\otimes\idon{\Qsys_s})$,
maps
\begin{equation}\label{eq:choi-feas}
  \Feas^{J}\coloneqq\Big\{\!(J_s)_{s\in\SecSet}:\ J_s\in\Pos{A_0\Qsys_s},\
  \sum_{s\in\SecSet}\tr_{\Qsys_s}J_s=\idon{A_0},\
  W_{\tau}\Big(\textstyle\sum_s
  \tr_{\Qsys_s}[(\idon{A_0}\otimes\hat O^{(s)}_k)J_s]\Big)
  W_{\tau}^{\dagger}\bowtie_k\Theta_k\ \forall k\!\Big\}
\end{equation}
\emph{onto} $\Feas$ of \eqref{eq:feas}, bijectively when the support frame
is tight (Lemma~\ref{lem:choi-marginal}), and, if $\Feas\ne\varnothing$ and
$\dim\Hs{\Qsys}\ge2$, only then. Writing
$\Phimap_{J}[J]\coloneqq\Phimap[\chi_{\Qsys}(J)]$,
$\cG^{J}_{\hat c,s}\coloneqq\cG_{\hat c,s}\circ\chi_{\Qsys_s}$,
$f^{J}_{\hat c,s}\coloneqq f_{\hat c,s}\circ\chi_{\Qsys_s}$,
$P^{J}_{\hat c}(J)\coloneqq\Phimap_J[J](\hat c)$ and
$K^{J}_{\hat c}(J)\coloneqq P^{J}_{\hat c}(J)
+\sum_{s\in\SecSet_K(\hat c)}f^{J}_{\hat c,s}(J_s)$, the value
\eqref{eq:hlow-def} equals
\begin{equation}\label{eq:choi-program}
\boxed{\ \
\begin{aligned}
  \hlow=\inf_{\mathbf q,\ (J_s)_s}\quad
  &\frac{\alpha}{\alpha-1}\sum_{\hat c\notin\Chkey}
  q(\hat c)\log_2\frac{q(\hat c)}{P^{J}_{\hat c}(J)}
  +\sum_{\hat c\in\Chkey}\left[
  q(\hat c)\log_2\frac{q(\hat c)}{P^{J}_{\hat c}(J)}
  +\frac{q(\hat c)}{\alpha-1}
  \log_2\frac{q(\hat c)}{K^{J}_{\hat c}(J)}\right]\\
  \text{s.t.}\quad
  &\mathbf q\in\acSet,\qquad (J_s)_s\in\Feas^{J},
\end{aligned}\ \ }
\end{equation}
an \emph{exact} reformulation (not a relaxation), jointly convex in
$(\mathbf q,(J_s)_s)$, with all constraints affine or conic.

\emph{In the general format} of Definition~\ref{def:declaration}, let the
frame be a direct sum $\iota=\bigoplus_{t\in\mathbb{T}}\iota_t$ with every
$\iota_t\in\Linn{A_{0,t}}{A_t}$ satisfying \eqref{eq:iota} for
$\tau^{(t)}_A$ --- either construction of Definition~\ref{def:support-frame}
applied sector by sector --- and put $W_t\coloneqq(\tau^{(t)}_A)^{1/2}\iota_t$. Let $\Feas^{J}$
be the set of lists $(J_{t,s})$ over the \tref[joint sector]{joint sectors} with
$J_{t,s}\in\Pos{A_{0,t}\Qsys_s}$, the coupling constraints
$\sum_s\tr_{\Qsys_s}J_{t,s}=\idon{A_{0,t}}$ for every $t$, and for every
entry the row of \eqref{eq:feas-general} evaluated at
$\omega^{(t,s)}=(W_t\otimes\idon{})J_{t,s}(W_t^{\dagger}\otimes\idon{})$: for
an entry read through its \tref{lifted operator},
$W_t\big(\sum_s\tr_{\Qsys_s}[(\idon{A_{0,t}}\otimes\hat O^{\uparrow(s)}_k)
J_{t,s}]\big)W_t^{\dagger}\bowtie_k\Theta^{(t)}_k$ for every $t$, and for a
\tref{party-local} entry the last line of \eqref{eq:feas-general} at $J$,
conjugated by $W_t\otimes\idon{}$. Then $(J_{t,s})\mapsto(\omega^{(t,s)})$
maps $\Feas^{J}$ onto $\Feas$ of \eqref{eq:feas-general}, bijectively when
every sector's frame is tight, and, if $\Feas\ne\varnothing$ and
$\dim\Hs{\Qsys}\ge2$, only then, and \eqref{eq:choi-program} holds with
$K^{J}_{\hat c}(J)=P^{J}_{\hat c}(J)+\sum_{(t,s)\in\mathbb{J}_K(\hat c)}
f_{\hat c,(t,s)}(\omega^{(t,s)})$. Every constraint of $\Feas^{J}$ involves
the blocks of one Alice sector only, so $\Feas^{J}=\prod_t\Feas^{J}_t$, and a
linear function of $(J_{t,s})$ is minimised over $\Feas^{J}$ by minimising,
independently for each $t$, its part on the blocks of $t$ over
$\Feas^{J}_t$.
\end{theorem}

\begin{proof}
\emph{Idea.} The map $\chi_{\Qsys}$ acts on the $A$ factor alone, so it
respects the sector blocks and carries every constraint row over by
substitution. Surjectivity is that of Lemma~\ref{lem:choi-marginal},
completed on one sector by a block that $W_{\tau}$ annihilates; the value
follows by substituting into the objective, and injectivity is not needed.

\emph{The map.} Lemma~\ref{lem:choi-marginal} restricts to
sector-block-diagonal operators because $\chi_{\Qsys}$ acts only on the
$A$ factor and therefore commutes with the sector projectors
$\idon{}\otimes\Pi_s$; the global trace-preservation condition
$\tr_{\Qsys}J=\idon{A_0}$ is the single coupling constraint
$\sum_s\tr_{\Qsys_s}J_s=\idon{A_0}$, which by the lemma \emph{implies}
$\omega_A=\tau_A$ and is implied by it up to the fibre the lemma
describes. The $\CList$ rows transfer by substitution, using
$\tr_{\Qsys}[(\idon{A}\otimes\hat O_k)\chi_{\Qsys}(J)]
=W_{\tau}(\tr_{\Qsys}[(\idon{A_0}\otimes\hat O_k)J])W_{\tau}^{\dagger}$
--- which holds because
$(\idon{A}\otimes\hat O_k)(W_{\tau}\otimes\idon{})
=(W_{\tau}\otimes\idon{})(\idon{A_0}\otimes\hat O_k)$ --- and
block-diagonality of $\hat O_k$: the $k$-th row of \eqref{eq:choi-feas}
\emph{is} the $k$-th row of \eqref{eq:feas-modes} evaluated at
$\omega=\chi_{\Qsys}(J)$, so it holds for $J$ exactly when it holds for
its image. Normalization $\tr\omega=1$ follows from the coupling
constraint and \eqref{eq:iota}, since
$\tr\omega=\tr[W_{\tau}^{\dagger}W_{\tau}\cdot\tr_{\Qsys}J]
=\tr[W_{\tau}^{\dagger}W_{\tau}]
=\tr[\iota^{\dagger}\tau_A\iota]=\tr[\tau_A\,\iota\iota^{\dagger}]
=\tr[\Pi_{\iota}\tau_A]=\tr\tau_A=1$. Together with positivity and
block-diagonality, which $\chi_{\Qsys_s}$ preserves, this gives
$\chi_{\Qsys}(\Feas^{J})\subseteq\Feas$.

\emph{Onto.} Fix $\omega\in\Feas$. Put
$J_{0,s}\coloneqq(W_{\tau}^{+}\otimes\idon{})\,\omega_s\,
(W_{\tau}^{+\dagger}\otimes\idon{})\succeq0$ for every $s$; then
$\chi_{\Qsys_s}(J_{0,s})=\omega_s$ exactly as in
Lemma~\ref{lem:choi-marginal} (each $\omega_s$ is supported on
$\operatorname{supp}\tau_A\otimes\Hs{\Qsys_s}$, which is
Lemma~\ref{lem:admissible-alice-subspace} with $\Sadm$ of \eqref{eq:SA}), and
$\sum_s\tr_{\Qsys_s}J_{0,s}
=W_{\tau}^{+}\big(\sum_s\tr_{\Qsys_s}\omega_s\big)W_{\tau}^{+\dagger}
=W_{\tau}^{+}\tau_AW_{\tau}^{+\dagger}=P_{\iota}$. Choose one
$s_0\in\SecSet$ ($\SecSet$ is non-empty, being the index set of a \tref{sector map},
\xref{foundations/def:sector-map}) and any
state $\varphi$ on $\Hs{\Qsys_{s_0}}$, and set
$J_s\coloneqq J_{0,s}+\delta_{s,s_0}\,(\idon{A_0}-P_{\iota})
\otimes\varphi$. Each $J_s\succeq0$;
$\sum_s\tr_{\Qsys_s}J_s=P_{\iota}+(\idon{A_0}-P_{\iota})=\idon{A_0}$; and
$W_{\tau}(\idon{A_0}-P_{\iota})=0$ kills the added block under both
$\chi_{\Qsys_{s_0}}$ and the conjugated constraint rows, so
$\chi_{\Qsys}(J)=\omega$ and every row of \eqref{eq:choi-feas} reads the
corresponding row of \eqref{eq:feas-modes} at $\omega$, which holds.
Hence $(J_s)_s\in\Feas^{J}$ and $\chi_{\Qsys}(\Feas^{J})=\Feas$.
Injectivity, when the frame is tight, is the lemma's. If the frame is not
tight, $\idon{A_0}-P_{\iota}\ne0$; if moreover $\Feas\ne\varnothing$ and
$\dim\Hs{\Qsys}\ge2$, there are two different choices of $(s_0,\varphi)$ ---
two sectors, or two states on one sector of dimension at least two --- and
for a fixed $\omega\in\Feas$ they give two different lists $(J_s)_s$ in
$\Feas^{J}$, differing in the added block, with the same image $\omega$.

\emph{The value.} $O^{J}_{\alpha}=O_{\alpha}\circ(\idmap\times
\chi_{\Qsys})$ by construction --- the objective transfers by
substitution of $\omega_s=\chi_{\Qsys_s}(J_s)$ into
\eqref{eq:objective}, using Definition~\ref{def:keyactive} for the
restriction of the $K$-sum to key-active sectors (an identity, by
Proposition~\ref{prop:exact-cancellation}) --- so
$\inf_{\acSet\times\Feas^{J}}O^{J}_{\alpha}
=\inf_{\mathbf q\in\acSet}\inf_{\omega\in\chi_{\Qsys}(\Feas^{J})}
O_{\alpha}(\mathbf q,\omega)
=\inf_{\mathbf q\in\acSet}\inf_{\omega\in\Feas}O_{\alpha}(\mathbf q,\omega)
=\hlow$, which is \eqref{eq:choi-program}. Only surjectivity enters this
chain and injectivity does not, so a non-tight support frame changes no
value.

\emph{Convexity.} Each $\chi_{\Qsys_s}$ is linear, so
Theorem~\ref{thm:convexity} applies composed with a linear map; the
constraint set \eqref{eq:choi-feas} is an intersection of the PSD cones
with affine conditions (for $\bowtie_k\in\{\preceq,\succeq\}$, a
conjugated affine-semidefinite condition), hence convex and compact ---
compact because $J_s\succeq0$ with
$\sum_s\tr J_s=\tr\idon{A_0}=d_{A_0}$.

\emph{The general format.} By Lemma~\ref{lem:choi-factor}\,(ii), applied with
the factors $W_t$ (factors of $\tau^{(t)}_A$ by \eqref{eq:WWtau} for each
sector), the assignment maps the lists satisfying the coupling constraints
onto the lists over the joint sectors with positive semidefinite blocks and
marginal $\tau_A$, injectively when every $\ker W_t=0$, which is
tightness of every sector's frame. Where some $\ker W_t\ne0$ and
$\dim\Hs{\Qsys}\ge2$, the onto step of that lemma's proof, run for $W_t$
with two different sector-block-diagonal states $\varphi\ne\varphi'$ (states
on two sectors, or two states on one sector of dimension at least two), gives
two different preimages of any point, differing by
$(\idon{A_{0,t}}-P_{W_t})\otimes(\varphi-\varphi')$; having the same
image, they satisfy the same rows, which $\Feas^{J}$ reads at the image
point, so at a point of a non-empty $\Feas$ both lie in $\Feas^{J}$.
Normalization follows from the marginal
condition and $\tr\tau_A=1$. The rows of $\Feas^{J}$ are by definition the
rows of \eqref{eq:feas-general} at the image point, so a list lies in
$\Feas^{J}$ exactly when it satisfies the coupling constraints and its image
lies in $\Feas$; for the lifted-operator rows the displayed conjugated form
is the substitution made in the paragraph \emph{The map}, sector by sector. Hence the image of
$\Feas^{J}$ is $\Feas$. The value: substitute into \eqref{eq:objective} and
use Lemma~\ref{lem:joint-sectors}\,(ii), (iii) for $K_{\hat c}$ and the
restriction of its sum to $\mathbb{J}_K(\hat c)$; surjectivity gives the
equality of the infima as before. Convexity and compactness are as before,
each coupling constraint bounding $\sum_s\tr J_{t,s}=d_{A_{0,t}}$. The
product structure: the coupling constraint of $t$, the lifted-operator rows
of $t$ and the party-local rows at $(t,s_{f^c})$ each involve only blocks
$J_{t,s}$ with that $t$, and a linear function on a Cartesian product is a
sum of linear functions of the factors, minimised factor by factor.
\end{proof}
